\section{The small-inclusion limit for an explicit parity reaction kernel}
\label{cl:section}

We now discharge the small-inclusion and whole-plane error estimates
for a specific limiting core Green kernel. This is a theorem about
that explicit kernel. Its identification as the limit of physical
finite radial backgrounds, together with control of the remote
annulus, is proved separately in Theorem~\ref{rg:physical-limit}.
Neither step is inferred just from finite-band incident amplification;
their combination gives Theorem~\ref{uc:counterexample}.

\subsection{The exact kernel and elementary plane estimates}

Write $k=k_0$, $\Phi_k(x)=e^{ik|x|}/(4\pi|x|)$, and use the
dimensionless electric polarization $P=\chi E$, where
$\chi=\eps_r-1$. The free electric response is
\begin{equation}\label{cl:free}
 \mathcal K_0P=(k^2I+\nabla\nabla)\Phi_k*P
       =-\mathcal P P+\mathcal D P.
\end{equation}
Here $\mathcal P$ is the longitudinal orthogonal projection on
$L^2(\mathbb R^3;\mathbb C^3)$, whose Fourier multiplier is
$\xi\xi^T/|\xi|^2$. Thus $\mathcal P\geq0$ and
$\|\mathcal P\|=1$. The remainder $\mathcal D$ has a kernel
bounded by $C/|x|$ on every fixed bounded set. Outside the support
of $P$ the static kernel $-\mathcal P$ is bounded by $C/|x|^3$.
The magnetic response is
\begin{equation}\label{cl:magnetic}
 \mathcal HP(x)=\frac{\nabla\times\mathcal K_0P(x)}{i\omega\mu_0}
 =\frac{k^2}{i\omega\mu_0}
       \int\nabla\Phi_k(x-y)\times P(y)\,dy.
\end{equation}
Its kernel is bounded by $C(|x|^{-2}+|x|^{-1})$.

Let $R_z=\operatorname{diag}(1,1,-1)$ and define the electric-field
wall projections
\begin{equation}\label{cl:projections}
 \mathcal W E(x)=-R_zE(R_zx),\qquad
 \Pi_{\rm e}=\tfrac12(I+\mathcal W),\qquad
 \Pi_{\rm m}=\tfrac12(I-\mathcal W).
\end{equation}
The first range has electric-wall parity. The entire real kernel
\begin{equation}\label{cl:J}
 J(x)=(k^2I+\nabla\nabla)\frac{\sin(k|x|)}{4\pi|x|}
 =\frac{k^3}{16\pi^2}\int_{S^2}
       (I-\theta\theta^T)e^{ik\theta\cdot x}\,d\theta
\end{equation}
is the imaginary part of the free electric kernel. Set
\[
 \mathcal R_{\rm e}P=i\Pi_{\rm e}(J*P),\qquad
 \mathcal R_{\rm m}P=i\Pi_{\rm m}(J*P).
\]
For an actual core gain parameter $B>0$, define
\begin{equation}\label{cl:kernel}
 \mathcal K_BP=\mathcal K_0P
       +(B^2-1)\mathcal R_{\rm e}P
       +(B^{-2}-1)\mathcal R_{\rm m}P.
\end{equation}
All derivatives in these definitions are distributional where
necessary. Away from the source support the formulas are classical.
Both reaction fields are entire vacuum Maxwell electric fields;
their magnetic partners are their curls divided by $i\omega\mu_0$.
On $z=0$ the electric-wall reaction has zero tangential electric
field, and the opposite-wall reaction has zero tangential magnetic
field.

\begin{lemma}[PROVED: uniform local and far-field kernel bounds]
\label{cl:kernelbounds}
Fix a ball containing all source supports under consideration.
On each fixed observation ball, $\mathcal R_{\rm e}P$,
$\mathcal R_{\rm m}P$, and their magnetic partners, with any fixed
number of derivatives, are bounded by $C\|P\|_1$.
For $|x|=r$ outside a larger fixed ball, each has an expansion
\begin{equation}\label{cl:transverse}
 F(x)=\frac{e^{ikr}F_+(\widehat x)+e^{-ikr}F_-(\widehat x)}r
                +O(C\|P\|_1r^{-2}),\qquad
 F_\pm(\widehat x)\cdot\widehat x=0,
\end{equation}
with $\|F_\pm\|_\infty\leq C\|P\|_1$. The free electric and
magnetic responses have the same statement with $F_-=0$.
For the dynamic electric response $\mathcal D P$ the same leading
far-field statement holds. These estimates are uniform over all
$L^1$ polarizations in that fixed source ball.

If two such fields have amplitude bounds $M_1,M_2$ in this
statement, then on the measurement plane outside the fixed larger
disk
\begin{equation}\label{cl:tailproduct}
 |(F_1\times\overline{F_2})\cdot\widehat z|
        \leq C M_1M_2r^{-3},\qquad
 \int_{r>R}|(F_1\times\overline{F_2})\cdot\widehat z|\,dA
        \leq C M_1M_2/R.
\end{equation}
The same assertions apply to any fixed finite collection of regular
vacuum waves, with its coefficient norm replacing $\|P\|_1$.
\end{lemma}
\begin{proof}
The integral in \eqref{cl:J} gives all compact derivative bounds
directly, including after reflection. The reaction fields are
smooth spherical Herglotz fields. Their angular amplitudes and all
their derivatives are bounded by $C\|P\|_1$, since $y$ stays in
a fixed ball in the factor $e^{-ik\theta\cdot y}$.
Stationary phase at the two points $\theta=\pm\widehat x$ gives
\eqref{cl:transverse}, uniformly in direction. The factor
$I-\theta\theta^T$ makes the electric leading amplitudes
transverse. A curl gives the transverse magnetic leading
amplitudes and the same remainder order. Equivalently these
statements follow by differentiating the explicit sine kernel.

For the free responses, expand $|x-y|=r-\widehat x\cdot y+O(r^{-1})$
and the derivatives of $\Phi_k(x-y)$ uniformly for bounded $y$.
The electric leading matrix is
$k^2(I-\widehat x\widehat x^T)e^{ikr}/(4\pi r)$;
the magnetic leading vector is a cross product with $\widehat x$.
Both are transverse. The static subtraction is $O(\|P\|_1r^{-3})$,
so it does not change the dynamic leading terms.
On the plane, $\widehat x\cdot\widehat z=0$. The cross product
of two vectors transverse to $\widehat x$ is parallel to
$\widehat x$, even with complex coefficients. Consequently every
$r^{-2}$ leading contribution to the normal cross product is zero.
The remainders give \eqref{cl:tailproduct}, whose radial integral
is finite. Finite regular waves have the same transverse incoming
and outgoing expansions.
\end{proof}

\begin{lemma}[PROVED: small-source plane estimates]
\label{cl:planebounds}
Let $\Omega_a$ be the union of the $2N$ radius-$\rho a\ell_0$
balls from Proposition~\ref{ct:model}, where $N$ is fixed and
$a\downarrow0$. They lie at distances at least
$(h-\rho)a\ell_0$ from the plane, and $|\Omega_a|=O(a^3)$.
For $P\in L^2(\Omega_a)$, extended by zero, and a fixed plane
disk $K$ containing all their projections,
\begin{align}
 \|\mathcal P P\|_{L^1(\mathbb R^2)}
       &\leq C a^{-1}\|P\|_1,\label{cl:staticL1}\\
 \|\mathcal P P\|_{L^2(\mathbb R^2)}
       &\leq C a^{-1/2}\|P\|_2,\label{cl:staticL2}\\
 \|\mathcal D P\|_{L^1(K)}
       &\leq C\|P\|_1,\label{cl:dynamicL1}\\
 \|\mathcal K_0P\|_{L^2(K)}
       &\leq C a^{-1/2}\|P\|_2,\qquad
 \|\mathcal HP\|_{L^2(K)}
       \leq C a^{1/2}\|P\|_2.\label{cl:nearL2}
\end{align}
The constants are independent of $a$ and $P$.
\end{lemma}
\begin{proof}
For a source point $y$ at height $d\geq c a$,
\[
 \int_{\mathbb R^2}|(x,0)-y|^{-3}\,dx=2\pi/d,
 \qquad
 \int_{\mathbb R^2}|(x,0)-y|^{-6}\,dx=\pi/(2d^4).
\]
The first proves \eqref{cl:staticL1}. The second, integrated over
$\Omega_a$, bounds the squared Hilbert--Schmidt norm by
$C a^3a^{-4}=C a^{-1}$, proving \eqref{cl:staticL2}.
The planar integral of $|x-y|^{-1}$ over $K$ is uniformly bounded,
which proves \eqref{cl:dynamicL1}.

For the magnetic squared kernel, the $|x-y|^{-4}$ planar
integral is $\pi/d^2$, and the $|x-y|^{-2}$ integral over $K$
is $O(1+|\log d|)$. Its squared Hilbert--Schmidt norm is thus
$O(a)+O(a^3|\log a|)=O(a)$. The dynamic electric squared
kernel has Hilbert--Schmidt norm squared $O(a^3|\log a|)$.
Combine these with the static bound to obtain \eqref{cl:nearL2}.
\end{proof}

\subsection{Self-consistent inclusions and the complete plane remainder}

Fix the interpolants $v_j$, their number $\lambda$, the physical
incoming basis $g_j$, and the geometry from
Lemmas~\ref{ct:jets}--\ref{ct:gram}. Let
$(E^v(c),H^v(c))$ be the finite regular electric-wall field
$\sum_jc_jv_j$. For fixed $t>0$ put
\begin{equation}\label{cl:scales}
 A_a=t/a,\qquad B_a=\sqrt\lambda A_a,\qquad
 (E^{\rm bg}_a,H^{\rm bg}_a)=A_a(E^v(c),H^v(c)),\qquad
 \chi_a=(\eps_p-1)\mathbf1_{\Omega_a}.
\end{equation}
The distinction between $A_a$ and $B_a$ is necessary: $A_a$ is
the jet scale, whereas $B_a$ is the actual common core multiplier
on the orthonormal physical inputs $g_j$.

\begin{theorem}[PROVED: explicit-kernel small-inclusion remainder]
\label{cl:remainder}
For all sufficiently small $a$, the equation
\begin{equation}\label{cl:integralequation}
 P_a=\chi_a\bigl(E^{\rm bg}_a+\mathcal K_{B_a}P_a\bigr)
 \quad\hbox{in }L^2(\Omega_a)
\end{equation}
has a unique solution. Define the auxiliary total fields by
\[
 E_a=E^{\rm bg}_a+\mathcal K_{B_a}P_a,
 \qquad H_a=\frac{\nabla\times E_a}{i\omega\mu_0}.
\]
They solve Maxwell's equations with the actual sphere permittivities
locally and outside the spheres, with the feedback condition encoded
by \eqref{cl:kernel}. If $F_a(c)$ is the explicit cross-profile
of Proposition~\ref{ct:model}, then
\begin{equation}\label{cl:actualremainder}
 \left\|\tfrac12\operatorname{Re}(E_a\times\overline{H_a})
                \cdot\widehat z-F_a(c)\right\|_{L^1(\mathbb R^2)}
       \leq C a\|c\|^2.
\end{equation}
All constants may depend on the fixed $N$, points, interpolants,
$\lambda,t,k,\rho,h,\eps_p$, but not on $a$ or $c$.
The theorem does not normalize these auxiliary fields by their own
free-space incoming amplitude. Their use for physical unit-power
scattering requires the separate finite-background limit.
\end{theorem}
\begin{proof}
It suffices to prove estimates for $\|c\|=1$; linearity of the
fields then gives the stated quadratic dependence. Write
$s_a=\sqrt{\chi_a}$ and $u_a=s_aE_a$, so $P_a=s_au_a$.
Equation \eqref{cl:integralequation} becomes
\begin{equation}\label{cl:positiveblock}
 \bigl(I+s_a\mathcal P s_a-s_a\mathcal D s_a
 -(B_a^2-1)s_a\mathcal R_{\rm e}s_a
 -(B_a^{-2}-1)s_a\mathcal R_{\rm m}s_a\bigr)u_a
      =s_aE^{\rm bg}_a.
\end{equation}
The first two terms form a self-adjoint operator at least $I$,
so its inverse has norm at most one. A Schur estimate for the
$|x-y|^{-1}$ dynamic kernel on $\Omega_a$ is $O(a^2)$:
within one ball integrate $|x-y|^{-1}$ over a ball of radius
$O(a)$, and between distinct balls use their fixed separation.
The smooth reaction kernels have operator norms $O(a^3)$.
Since $B_a=O(a^{-1})$, the norm of the remaining perturbation
in \eqref{cl:positiveblock} is $O(a^2)+O(B_a^2a^3)=O(a)$.
A Neumann series after the positive inverse proves invertibility,
with inverse norm at most two for small $a$. In particular
\begin{equation}\label{cl:sourcebounds}
 \|P_a\|_2\leq C a^{1/2},\qquad \|P_a\|_1\leq C a^2.
\end{equation}
The equation and all its kernels commute with the rotation by
$\pi$ about $z$; the source and solution are in its invariant
sector. No unsupported odd-$m$ extension of a core kernel is used.

Let $P_a^0$ be constant on each sphere, with the isolated static
value
\begin{equation}\label{cl:leadingP}
 P_a^0\big|_{\Omega_{a,j}^{\pm}}
   =(\eps_p-1)\frac3{\eps_p+2}
       E^{\rm bg}_a(p_j^{\pm}),\qquad p_j^+=p_j, p_j^-=-p_j.
\end{equation}
These vectors are normal to the plane. For one isolated sphere
they solve the static equation exactly: the interior electric
field of its constant polarization is $-P_a^0/3$, so
$P_a^0=\chi(E^{\rm bg}_a-P_a^0/3)$.
The actual background differs from its value at $p_j^{\pm}$
by $O(A_aa)=O(1)$ throughout the corresponding sphere. Its
$L^2$ residual is $O(a^{3/2})$.
Static interaction between distinct spheres has operator norm
$O(a^3)$, because their mutual distances remain positive.
The dynamic and reaction perturbations on $P_a^0$ have norm
$O(a)\|P_a^0\|_2=O(a^{3/2})$ by the preceding bounds.
Applying the uniformly bounded inverse in \eqref{cl:positiveblock}
to this residual proves
\begin{equation}\label{cl:sourceerror}
 \|P_a-P_a^0\|_2\leq C a^{3/2},\qquad
 \|P_a-P_a^0\|_1\leq C a^3.
\end{equation}

It remains to estimate the whole plane, not merely a compact set.
Abbreviate the free response by $(E^f,H^f)$, and the two full
reaction responses, including their coefficients in
\eqref{cl:kernel}, by $(E^e,H^e)$ and $(E^m,H^m)$.
On the plane the wall identities give exactly
\begin{equation}\label{cl:fluxsplit}
 E_{a,t}=E^f_t+E^m_t,\qquad
 H_{a,t}=H^{\rm bg}_{a,t}+H^f_t+H^e_t.
\end{equation}
By Lemma~\ref{cl:kernelbounds} and \eqref{cl:sourcebounds}, the
opposite-wall reaction fields have compact and far-field amplitude
$O(a^2)$, while the wall reaction fields have amplitude
$O(B_a^2a^2)=O(1)$. The background fields have amplitude $O(A_a)$.

First take the cross product with $H^{\rm bg}_a$ in
\eqref{cl:fluxsplit}. Replacing the static free electric response
of $P_a$ by that of $P_a^0$ has plane $L^1$ error at most
$C a^{-1}\|P_a-P_a^0\|_1=O(a^2)$ by
\eqref{cl:staticL1}. Since the background is globally bounded by
$C A_a$, its flux error is $O(a)$. The dynamic free electric
response has $L^1$ norm $O(a^2)$ on a fixed disk, by
\eqref{cl:dynamicL1}; outside that disk the transverse cancellation
\eqref{cl:tailproduct} gives a flux bound
$O(A_a\|P_a\|_1)=O(a)$. The same compact bound and tail argument
apply to $E^m\times\overline{H^{\rm bg}_a}$.

The remaining static field of $P_a^0$ is exactly a sum of normal
dipoles, with the polarizabilities from Lemma~\ref{ct:profile}.
For each individual dipole, its tangential plane trace satisfies
\begin{equation}\label{cl:firstmoment}
 \int_{\mathbb R^2}|E^{0,j}_t(x)|\,dA\leq C a,
 \qquad
 \int_{\mathbb R^2}|x-p_j^{\pm}|\,|E^{0,j}_t(x)|\,dA
          \leq C a^2.
\end{equation}
These follow directly by scaling \eqref{ct:dipoletrace}; the
second radial integrand at infinity is $O(r^{-2})$ and is
integrable. A finite regular background has global Lipschitz
constant $C A_a$. Replacing its magnetic field in each such cross
product by $H^{\rm bg}_a(p_j^{\pm})$ therefore costs at most
$C A_a a^2=O(a)$. The resulting sum is exactly $F_a(c)$, including
the coefficients and signs in \eqref{ct:profileformula}.

We finally control all cross products with the two correction
magnetic fields in \eqref{cl:fluxsplit}. On a fixed disk,
\eqref{cl:nearL2} and \eqref{cl:sourcebounds} give
\[
 \|E^f\|_2\leq C,
 \qquad \|H^f\|_2\leq Ca,
 \qquad \|E^m\|_\infty\leq Ca^2,
 \qquad \|H^e\|_\infty\leq C.
\]
Thus the free self-cross product and the product
$E^m\times\overline{H^f}$ have $L^1$ norm $O(a)$.
For products with $H^e$, use
$\|\mathcal P P_a\|_{L^1}\leq C a^{-1}\|P_a\|_1=O(a)$,
$\|\mathcal D P_a\|_{L^1(K)}=O(a^2)$, and the bounded
opposite-wall reaction, obtaining another $O(a)$ bound.
Outside the disk, \eqref{cl:tailproduct} applies to these products
with amplitude pairs $(a^2,a^2)$ or $(a^2,1)$; both give $O(a)$
or smaller. All normal products are therefore integrable, and
their sum proves \eqref{cl:actualremainder}.
\end{proof}

\begin{remark}[What has been discharged]
\label{cl:scope}
For the explicit kernel \eqref{cl:kernel}, the sufficient plane
remainder in Proposition~\ref{ct:completion} is now a conclusion,
not an assumption. The proof uses the uniformly positive static
block, includes every coherent input, and estimates the full
absolute plane flux. Theorem~\ref{rg:physical-limit} supplies the
separate physical sequence: its core source Green kernels converge
to \eqref{cl:kernel}, and its remote-annulus and exterior plane-flux
contributions tend to zero at each fixed $a$.
The order is fixed $a$ first, remote radius second, then a diagonal
sequence with $a\downarrow0$. Theorem~\ref{uc:counterexample}
combines these independently proved limits into finite ordinary
devices with the original physical incoming normalization.
\end{remark}
